\chapter{Soluzioni agli esercizi}

In questo capitolo sono riportate le soluzioni agli esercizi proposti nei vari capitoli del libro. E' chiaro che ci possano essere dei modi alternativi per risolvere gli esercizi, ma le soluzioni che propongo qua sono quelle che mi sono sembrate più semplici e dirette.

\section{Spazi vettoriali}
\subsection*{Esercizio \ref{ex:1.1}}
Per quanto riguarda la dimostrazione che il vettore nullo è unico, possiamo procedere per assurdo: negando la tesi, si ha che esistono due elementi distinti $\underline{0}, \underline{0}' \in V$ tali che
    $$
    \forall \underline{v} \in V, \underline{0} + \underline{v} = \underline{v} + \underline{0} = \underline{v} = \underline{0}' + \underline{v} = \underline{v} + \underline{0}'.
    $$
    Ma, osservando che queste due relazioni devono valere per ogni $\underline{v}$, si ha che
    \begin{align*}
\underline{0}' &= \underline{0}' + \underline{0} & \text{(perché \underline{0} è l'elemento neutro)} \\
&= \underline{0} + \underline{0}' & \text{(perché \underline{0}' è un elemento neutro rispetto a $+$)} \\
&= \underline{0} & \text{(perché \underline{0} è un elemento neutro rispetto a $+$)}.
\end{align*}
Ma questo è una contraddizione, in quanto abbiamo assunto che $\underline{0}$ e $\underline{0}'$ fossero distinti. Segue la tesi. \\
Per quanto riguarda la dimostrazione che $\forall \underline{v} \in V$ il vettore nullo è $0 \cdot \underline{v}$, si osserva che
\begin{align*}
    0_\mathbb{K} \cdot \underline{v} = (0_\mathbb{K} +_{\mathbb{K}} 0_\mathbb{K}) \cdot \underline{v} = 0_{\mathbb{K}} \cdot \underline{v} + 0_\mathbb{K} \cdot \underline{v}, 
    \end{align*}
dove si è usato il fatto che $0_\mathbb{K} = 0_\mathbb{K} +_{\mathbb{K}} 0_\mathbb{K}$ (essendo l'elemento neutro rispetto all'operazione $+_\mathbb{K}$ definita sul campo $\mathbb{K}$) e la proprietà distributiva di cui gode $\cdot$ rispetto a $+_\mathbb{K}$. A questo punto, sappiamo che $\forall \underline{v} \in V, \exists (-\underline{v}) \in V$ tale che $\underline{v} + (-\underline{v}) = (-\underline{v}) + \underline{v} = \underline{0}$, dunque sappiamo che esiste un elemento inverso per $0_{\mathbb{K}} \cdot \underline{v}$ che possiamo aggiungere a destra di ogni membro\footnote{esplicito questa cosa di aggiungere a sinistra perché, a priori, queste dimostrazioni si potrebbero fare anche su spazi dove non è garantita la proprietà commutativa della somma, ottenendo dei risultati simili, ma in questo caso non sarebbe necessario. La faccio per pura pedanteria e perché alcune di queste dimostrazioni potrebbero esseri riviste a Teoria dei Gruppi}
    $$
    0_\mathbb{K} \cdot \underline{v} + (-0_\mathbb{K} \cdot \underline{v}) = 0_\mathbb{K} \cdot \underline{v} + 0_\mathbb{K} \cdot \underline{v} + (-0_\mathbb{K} \cdot \underline{v}).
    $$
Ma, essendo $0_\mathbb{K} \cdot \underline{v} + (-0_\mathbb{K} \cdot \underline{v}) = \underline{0}$, si ha che
    $$
    \underline{0} = 0_\mathbb{K} \cdot \underline{v} + \underline{0} = 0_\mathbb{K} \cdot \underline{v},
    $$
dove si è usato il fatto che $\underline{0}$ è l'elemento neutro rispetto all'operazione $+$ definita su $V$. Segue la tesi. \\
Per quanto riguarda l'ultima proposizione da mostrare abbiamo, per quanto visto nella proposizione \ref{prop:0v_equal_0}, che
$$
    \underline{0} = 0_\mathbb{K} \cdot \underline{v} = (1_\mathbb{K} +_\mathbb{K} (-1_\mathbb{K})) \cdot \underline{v} = \underline{v} + (-1_\mathbb{K}) \cdot \underline{v},
$$
dove si è usata la proprietà distributiva di cui gode l'operazione $\cdot$ rispetto all'operazione $+_\mathbb{K}$. A questo punto, sappiamo che $\forall \underline{w} \in V, \exists (-\underline{w}) \in V$ tale che $\underline{w} + (-\underline{w}) = (-\underline{w}) + \underline{w} = \underline{0}$, dunque sappiamo che esiste un elemento inverso per $\underline{v}$ che denoteremo sempre con $(-\underline{v})$ e che possiamo aggiungere a sinistra di ogni membro
$$
    (- \underline{v}) + \underline{0} = (-\underline{v}) + \underline{v} + (-1_\mathbb{K}) \cdot \underline{v} \implies (-\underline{v}) = \underline{0} + (-1_\mathbb{K}) \cdot \underline{v} = (-1_\mathbb{K}) \cdot \underline{v},
$$
dove abbiamo sempre usato il fatto che $\underline{0}$ è l'elemento neutro rispetto all'operazione $+$ definita su $V$. Segue la tesi

\subsection*{Esercizio \ref{ex:1.2}}
Per quanto riguarda la dimostrazione che $\text{Span}(\underline{v_1}, \ldots, \underline{v_n})$ è un sottospazio vettoriale, osserviamo che possiamo usare la proposizione \ref{prop:ssp_vett_iff}, dunque è sufficiente mostrare che $\forall \underline{v}, \underline{w} \in V$ e $\forall \alpha \in \mathbb{K}$ valgono le implicazioni $\underline{v}, \underline{w} \in W \implies \underline{v}+\underline{w} \in W$ e $\underline{v} \in V \implies \alpha \underline{v} \in V$ (dove ho sottointeso il prodotto per scalare). \\
Per quanto riguarda la prima implicazione, è chiaro che se $\underline{v},\underline{w} \in V$, allora per definizione esistono $\alpha_1, \ldots, \alpha_n, \beta_1, \ldots, \beta_n \in \mathbb{K}$ tali che $\underline{v} = \alpha_1 \underline{v_1} + \ldots + \alpha_n \underline{v_n}$ e $\underline{w} = \beta_1 \underline{v_1} + \ldots + \beta_n \underline{v_n}$, dunque se sommiamo i due vettori otteniamo che
\begin{align*}
    &\underline{v} + \underline{w} = (\alpha_1 +_\mathbb{K} \beta_1) \underline{v_1} + \ldots + (\alpha_n +_\mathbb{K} \beta_n) \underline{v_n} \implies \exists \gamma_1 = \alpha_1 + \beta_1, \ldots, \gamma_n = \alpha_n + \beta_n \in \mathbb{K} \text{ tali che } \\
    &\underline{v} + \underline{w} = \gamma_1 \underline{v_1} + \ldots + \gamma_n \underline{v_n} \implies \underline{v} + \underline{w} \in W.
\end{align*}
Per quanto riguarda la seconda implicazione, è chiaro che se $\underline{v} \in V$, allora per definizione esistono $\alpha_1, \ldots, \alpha_n \in \mathbb{K}$ tali che $\underline{v} = \alpha_1 \underline{v_1} + \ldots + \alpha_n \underline{v_n}$, dunque
\begin{align*}
    &\alpha \underline{v} = \alpha (\alpha_1 \underline{v_1} +_\mathbb{K} \ldots +_\mathbb{K} \alpha_n \underline{v_n}) = (\alpha \alpha_1) \underline{v_1} + \ldots + (\alpha \alpha_n) \underline{v_n} \implies \exists \beta_1 = \alpha \alpha_1, \ldots, \beta_n = \alpha \alpha_n \in \mathbb{K} \\
    &\text{ tali che } \alpha \underline{v} = \beta_1 \underline{v_1} + \ldots + \beta_n \underline{v_n} \implies \alpha \underline{v} \in W.
\end{align*}
Si osservi che in entrambe le implicazioni abbiamo usato il fatto che le operazioni $+$ e $\cdot$ godessero delle proprietà distributive rispetto a $+_{\mathbb{K}}$. Segue la tesi. \\
Per quanto riguarda la dimostrazione che $\text{Span}(\underline{v_1}, \ldots, \underline{v_n})$ è il più piccolo sottospazio vettoriale contenente $\{\underline{v_1}, \ldots, \underline{v_n}\}$, osserviamo che, per definizione di spazio vettoriale, ogni sottospazio vettoriale $W$ di $V$ che contiene gli elementi $\underline{v_1}, \ldots, \underline{v_n}$, contiene anche le combinazioni lineari di questi elementi, ovvero contiene $\text{Span}(\underline{v_1}, \ldots, \underline{v_n})$. Segue la tesi.
\subsection*{Esercizio \ref{ex:1.3}}
Per quanto riguarda la dimostrazione che $\text{Span}(A)$ è un sottospazio vettoriale, osserviamo che possiamo sempre usare la proposizione \ref{prop:ssp_vett_iff}, dunque è sufficiente mostrare che valgono le due implicazioni come fatto nell'esercizio precedente. \\
Osserviamo che preso $\underline{v} \in \text{Span}(A)$ allora sappiamo che $\exists \underline{v_1}, \ldots, \underline{v_n} \in A : \exists \alpha_1, \ldots, \alpha_n \in \mathbb{K} : \underline{v} = \alpha_1 \underline{v_1} + \ldots + \alpha_n \underline{v_n}$ e, similmente, se $\underline{w} \in \text{Span}(A)$ allora $\exists \underline{w_1}, \ldots, \underline{w_m} \in A : \exists \beta_1, \ldots, \beta_m \in \mathbb{K}$; dunque sommandoli insieme
\begin{align*}
    &\underline{v} + \underline{w} = (\sum_{i=1}^n \alpha_i \underline{v_i}) + (\sum_{i=1}^m \beta_i \underline{w_i}) = \alpha_1 \underline{v_1} + \ldots + \alpha_n \underline{v_n} + \beta_1 \underline{w_1} + \ldots + \beta_m \underline{w_m} \implies \\
    &\exists \underline{v_1}, \ldots, \underline{v_n}, \underline{w_1}, \ldots, \underline{w_m} \in A : \exists \alpha_1, \ldots, \alpha_n, \beta_1, \ldots, \beta_m \in \mathbb{K} : \underline{v} = \sum_{i=1}^n \alpha_i \underline{v_i} + \sum_{i=1}^n \beta_i \underline{w_i} \implies \\
    &\implies \underline{v} + \underline{w} \in \text{Span}(A).
\end{align*}
Mostriamo la seconda implicazione: come prima, se $\underline{v} \in \text{Span}(A)$ allora $\exists \underline{v_1}, \ldots, \underline{v_n} \in A : \exists \alpha_1, \ldots, \alpha_n \in \mathbb{K} : \underline{v} = \alpha_1 \underline{v_1} + \ldots + \alpha_n \underline{v_n}$, dunque preso $\alpha \in \mathbb{K}$ abbiamo che
\begin{align*}
&\alpha \underline{v} = \alpha (\alpha_1 \underline{v_1} + \ldots + \alpha_n \underline{v_n}) = (\alpha \alpha_1) \underline{v_1} + \ldots + (\alpha \alpha_n) \underline{v_n} \implies \\
&\implies \exists \underline{v_1}, \ldots, \underline{v_n} \in A : \exists \beta_1 = \alpha \alpha_1, \ldots, \beta_n = \alpha \alpha_n \in \mathbb{K} : \alpha \underline{v} = \sum_{i=1}^n \beta_i \underline{v_i} \implies \alpha \underline{v} \in \text{Span}(A).
\end{align*}
Si osservi che in entrambe le implicazioni abbiamo usato il fatto che le operazioni $+$ e $\cdot$ godano delle proprietà distributive rispetto a $+_\mathbb{K}$. Segue la tesi. \\
Per quanto riguarda il dimostrare che se $A$ è un sottospazio vettoriale allora $A = \text{Span}(A)$, è sufficiente mostrare che $A \subseteq \text{Span}(A)$ (che è banale) e che $\text{Span}(A) \subseteq A$, da cui segue che $A = \text{Span}(A)$. \\
$\boxed{\subseteq}$: preso $\underline{v} \in A$, allora per definizione di $\text{Span}(A)$ dobbiamo mostrare $\exists \underline{v_1}, \underline{v_2}, \ldots, \underline{v_n} \in A : \exists \alpha_1, \ldots, \alpha_n \in \mathbb{K} : \underline{v} = \alpha_1 \underline{v_1} + \ldots + \alpha_n \underline{v_n}$, ma possiamo direttamente prendere $\underline{v_1} = \underline{v}$ e $\alpha_1 = 1_\mathbb{K}$ (la definizione di $\text{Span}$ non ci vieta di farlo). Dunque $\underline{v} \in \text{Span}(A)$ e, per l'arbitrarietà di $\underline{v}$, segue che $A \subseteq \text{Span}(A)$. \\
$\boxed{\supseteq}$: preso $\underline{v} \in \text{Span}(A)$, abbiamo per definizione di $\text{Span}(A)$ che $\exists \underline{v_1}, \ldots, \underline{v_n} \in A : \exists \alpha_1, \ldots, \alpha_n \in \mathbb{K} : \underline{v} = \alpha_1 \underline{v_1} + \ldots + \alpha_n \underline{v_n}$ ma $\alpha_1 \underline{v_1} + \ldots + \alpha_n \underline{v_n} \in A$ (essendo un sottospazio vettoriale), dunque $\underline{v} \in A$ e, per l'arbitrarietà di $\underline{v}$, segue che $\text{Span}(A) \subseteq A$. \\
Segue allora che $A = \text{Span}(A)$, ovvero la tesi. \\
Mostriamo adesso che se $A \subseteq \text{Span}(B) \implies \text{Span}(A) \subseteq \text{Span}(B)$: è chiaro che se $A \subseteq \text{Span}(B)$, allora $\forall \underline{v} \in A, \exists \underline{v_1}, \ldots, \underline{v_n} \in B : \exists \alpha_1, \ldots, \alpha_n \in \mathbb{K} : \underline{v} = \alpha_1 \underline{v_1} + \ldots + \alpha_n \underline{v_n}$, dunque se consideriamo un vettore $\underline{w} \in \text{Span(A)}$ abbiamo per definizione che $\exists \underline{v_1'}, \ldots, \underline{v_n'} \in A : \exists \alpha_1', \ldots, \alpha_n' \in \mathbb{K}$ tali che
$$
    \underline{w} = \alpha_1' \underline{v_1'} + \ldots + \alpha_n' \underline{v_n'}.
$$
Ma, per quanto detto sopra, si ha che $\forall \underline{v_i'}, \exists \underline{v^{(i)}_1}, \ldots, \underline{v^{(i)}_m} \in A : \exists \beta^{(i)}_1, \ldots, \beta^{(i)}_n \in \mathbb{K} : \underline{v_i'} = \beta^{(i)}_1 \underline{v^{(i)}_1} + \ldots + \beta^{(i)}_n \underline{v^{(i)}_n}$, dunque sostituendo nella definizione di $\underline{w}$ si ottiene una combinazione lineare di vettori di $B$, dunque $\underline{w} \in \text{Span}(B)$. Segue la tesi. \\
Per quanto riguarda l'ultima proposizione da mostrare, è sufficiente ricondursi a mostrare che se $A \subseteq B \implies A \subseteq \text{Span}(B)$ e, per quanto visto nella proposizione precedente, ciò implica che $\text{Span}(A) \subseteq \text{Span}(B)$. Ma è chiaro che $B \subseteq \text{Span}(B)$, dunque $A \subseteq B \subseteq \text{Span}(B) \implies A \subseteq \text{Span}(B) \implies \text{Span}(A) \subseteq \text{Span}(B)$.
\subsection*{Esercizio \ref{ex:1.4}}

Possiamo far vedere che presa una collezione finita di polinomi non è possibile generare tutto lo spazio: consideriamo i polinomi $p_1(x), \ldots, p_n(x)$ e consideriamo lo Span di $p_1, \ldots, p_n$. Osserviamo che $\text{deg}(\alpha_1 p_1(x) + \ldots + \alpha_n p_n(x)) \leq \max (\text{deg}(p_1), \text{deg}(p_2), \ldots, \text{deg}(p_n))$ dove $\text{deg}$ indica il grado del polinomio; tuttavia è chiaro che $\mathbb{R}[x] \ni x^{1 + \max (\text{deg}(p_1), \text{deg}(p_2), \ldots, \text{deg}(p_n))} \not\in \text{Span}(p_1, \ldots, p_n)$.

\subsection*{Esercizio \ref{ex:1.5}}
Supponiamo per assurdo che $\exists A' \subseteq A: A'$ non è linearmente indipendente, allora sappiamo che $\exists \underline{v_1}, \ldots, \underline{v_n} \in A': \exists \alpha_1, \ldots, \alpha_n \in \mathbb{K}$ non tutti uguali a $0_\mathbb{K}$ tali che
$$
\alpha_1 \underline{v_1} + \ldots + \alpha_n \underline{v_n} = \underline{0},
$$ 
ma questa si tratta anche di una relazione di vettori appartenenti anche ad $A$, quindi $A$ sarebbe linearemente dipendente ottenendo così un assurdo.

\subsection*{Esercizio \ref{ex:1.6}}
Preso $\underline{v} \in A' \setminus A$ abbiamo che, per ipotesi, $\underline{v} \in \text{Span}(A)$ (siccome $A$ insieme di generatori), dunque $A' \subseteq \text{Span}(A) \implies \text{Span}(A') \subseteq \text{Span}(A) = V$, ma è chiaro che $A \subseteq A' \implies \text{Span}(A) \subseteq \text{Span}(A')$ (dove si è usato, nei due casi, la proposizione \ref{prop:incl_sub} e \ref{prop:incl_span} rispettivamente). Concludiamo allora che $\text{Span}(A) = \text{Span}(A') = V$, dunque anche $A'$ è un insieme di generatori.

\subsection*{Esercizio \ref{ex:1.7}}
Supponiamo che $A = \{ \underline{v_1}, \ldots, \underline{v_n} \}$ insieme di generatori e osserviamo che possiamo sempre procedere come fatto nella dimostrazione dell'esistenza di una base, ricavandoci un insieme linearmente indipendente massimale: se $A$ è linearmente indipendente, chiaramente, non c'è niente da fare. Se $A$ invece non lo è possiamo considerare la relazione non banale
$$
    \alpha_1 \underline{v_1} + \ldots + \alpha_n \underline{v_n} = \underline{0}
$$
e supponendo, senza perdita di generalità, che $\alpha_n \neq 0$ allora
$$
    \underline{v_n} = - \frac{\alpha_1}{\alpha_n} \underline{v_1} - \ldots - \frac{\alpha_{n-1}}{\alpha_n} \underline{v_{n-1}},
$$
quindi si ha, in virtù del lemma \ref{lemma:rim_superf}, che $\text{Span}(A) = \text{Span}(A \setminus \{ \underline{v_n} \})$, quindi possiamo limitarci a considerare l'insieme $A \setminus \{ \underline{v_n} \}$ e iterare il ragionamento finché non si ottiene un insieme linearmente indipendente: questo lemma ci assicura che lo Span rimane invariato.

\subsection*{Esercizio \ref{ex:1.8}}
Se considero un insieme di vettori linearmente indipendenti $B$ è chiaro che se generano $V$ non c'è niente da fare siccome rappresentano già una base dello spazio vettoriale. Se così non è, abbiamo mostrato che è sempre possibile trovare una base $\mathcal{B}$ nel nostro spazio vettoriale e, usando il lemma \ref{lemma:alg_scambio_pot}, sappiamo che è possibile trovare $\mathcal{B'} \subseteq \mathcal{B}$ tale che $(\mathcal{B} \setminus \mathcal{B'}) \cup B$ è ancora linearmente indipendente e ha lo stesso $\text{Span}$ dell'insieme iniziale, dunque è effettivamente una base del nostro spazio vettoriale.

\subsection*{Esercizio \ref{ex:1.9}}
Sappiamo che $W$ essendo sottospazio vettoriale, quindi pur sempre uno spazio vettoriale, ammette una base $\mathcal{W} = \{ \underline{w_1}, \ldots, \underline{w_m} \}$: tale insieme rimane comunque un insieme linearmente indipendente di vettori di $V$, pertanto possiamo applicare il lemma \ref{lemma:alg_scambio_pot} ad una base $\mathcal{B}$ di $V$ (che sappiamo esistere) per dedurre che $\# \mathcal{W} \leq \# \mathcal{B}$, ovvero $\text{dim}(W) \leq \text{dim}(V)$. Per quanto riguarda l'ultima richiesta dell'esercizio, osserviamo che l'implicazione verso sinistra $\boxed{\Leftarrow}$ è ovvia: se $V = W$ allora le dimensioni di $V$ e $W$ coincidono siccome una base di uno è base anche dell'altro; ma è altrettanto banale l'implicazione verso destra $\boxed{\Rightarrow}$ siccome si vede che se $\text{dim}(W) = \text{dim}(V)$ e $W \subseteq V$ allora una base di $n = \text{dim}(W)$ vettori di $W$ è anche una base di $V$ (siccome si parla di $n = \text{dim}(W) = \text{dim}(V)$ vettori linearmente indipendenti di $W$ e, dunque, di $V$).

\subsection*{Esercizio \ref{ex:1.10}}
Siccome $Z \subseteq V$ è un sottospazio vettoriale, dunque uno spazio vettoriale allora esiste una base $\mathcal{B} = \{ \underline{v_1}, \ldots, \underline{v_k} \}$ che possiamo completare a base di $\mathcal{B'} = \{ \underline{v_1}, \ldots, \underline{v_k}, \underline{v_{k+1}}, \ldots, \underline{v_n} \}$, dunque $k \leq n$. Se $k = n$ allora estendendo vettori non aggiungo nulla, quindi $V = W$.

\subsection*{Esercizio \ref{ex:1.11}}
Supponiamo per assurdo che $U \cap W \neq \{ \underline{0} \}$: preso $\underline{v} \in U \cap W$, si ha, per definizione di intersezione, che $\underline{v} \in U$ e $\underline{v} \in W$: preso a questo punto un vettore $\underline{z} \in Z$ sappiamo per ipotesi che $\exists! \underline{u'} \in U, \underline{w'} \in W$ tali che $\underline{z} = \underline{u'} + \underline{w'}$ ma
$$
    \underline{z} = \underline{u'} + \underline{w'} = \underline{u'} + \underline{v} - \underline{v} + \underline{w'} = \underset{\in U}{(\underline{u'} + \underline{v})} + \underset{\in W}{(\underline{w'} - \underline{v})},
$$
dove abbiamo il fatto che la legge di composizione interna produce vettori appartenenti sempre al sottospazio considerato. Segue allora che ogni vettore non può più essere scritto in maniera univoca, ottenendo così un assurdo. Segue così la tesi.
\subsection*{Esercizio \ref{ex:1.12}}
Supponiamo, senza perdita di generalità, che $A' = \{ \underline{v_1}, \ldots, \underline{v_k} \}$ e $A'' = \{ \underline{v_{k+1}}, \ldots, \underline{v_n} \}$ (a meno di riarrangiare i termini), osserviamo che preso $\underline{v} \in \text{Span}(A)$ allora $\exists! \alpha_1, \ldots, \alpha_n \in \mathbb{K}$ tali che
\begin{align*}
    &\underline{v} = \alpha_1 \underline{v_1} + \ldots + \alpha_n \underline{v_n} = \underset{\in \text{Span}(A')}{(\alpha_1 \underline{v_1} + \ldots + \alpha_k \underline{v_k})} + \underset{\in \text{Span}(A'')}{(\alpha_{k+1} \underline{v_{k+1}} + \ldots + \alpha_n \underline{v_n})} \implies \\
    &\implies \text{Span}(A) = \text{Span}(A') + \text{Span}(A'').
\end{align*}
Rimane da mostrare che $\text{Span}(A') \cap \text{Span}(A'') = \{ \underline{0} \}$: se per assurdo non fosse così allora preso $\underline{v} \in \text{Span}(A') \cap \text{Span}(A'')$ sappiamo che $\underline{v} \in \text{Span}(A')$ e $\underline{v} \in \text{Span}(A'')$, dunque $\exists \alpha_1, \ldots, \alpha_n \in \mathbb{K}$ non tutti nulli tali che
$$
    \underline{v} = \alpha_1 \underline{v_1} + \ldots + \alpha_k \underline{v_n} = \alpha_{k+1} \underline{v_{k+1}} + \ldots + \alpha_n \underline{v_n} \implies \alpha_1 \underline{v_1} + \ldots - \alpha_n \underline{v_n} = \underline{0},
$$
che rappresenta una relazione non banale fra vettori $\underline{v_1}, \ldots, \underline{v_n}$ che si erano supposti linearmente indipendenti, ovvero un assurdo!

\section*{Esercizio 2.*}
Osserviamo che la dimensione dello spazio delle matrici simmetriche $U$ è pari a $\frac{n(n+1)}{2}$ e la dimensione delle spazio delle matrici antisimmetriche $Z$ è pari a $\frac{n(n-1)}{2}$. E' chiaro che $U \cap Z = \{ \emptyset \}$ e vediamo che
$$
    \text{dim}(U) + \text{dim}(Z) = \frac{n(n+1)}{2} + \frac{n(n-1)}{2} = \frac{n^2 + n + n^2 - n}{2} = \frac{2n^2}{2} = n^2 = \text{dim}(\mathcal{M}_n(\mathbb{R})),
$$
e da qua possiamo concludere. Possiamo trovare esplicitamente tali matrici: abbiamo infatti che $U = A + A^t$ è simmetrica, siccome $U^t = (A + A^t)^t = A^t + (A^t)^t = A^t + A = A + A^t = U$ mentre la matrice $Z = A - A^t$ è antisimmetrica, siccome $Z^t = (A - A^t)^t = A^t - (A^t)^t = A^t - A = - (A - A^t) = - Z$ e
$$
    U + Z = A + A^t + A - A^t = 2A \implies A = \frac{U}{2} + \frac{Z}{2} = \frac{1}{2}(A+A^t) + \frac{1}{2}(A - A^t).
$$